\section{Motivation}
\label{section:motivation}

\subsection{Computational Interpretation of SAX}

Beyond its proof-theoretic results, SAX admits a natural operational
interpretation as a system of concurrent processes that interact
asynchronously. This interpretation has been studied across several
computational models, from shared memory concurrency \cite{DeYoung2020}
to message-passing concurrency \cite{}, and each of these interpretations has
been shown to satisfy desirable properties such as progress, preservation,
confluence, and termination. Specifically, under the shared memory
interpretation, variables are substituted by addresses at runtime, and thus the
sequent
\[
  x_1 : A_1, \ldots, x_n : A_n \vdash P :: (z : C)
\]
should be understood as typing a process $P$ that reads from addresses $x_1,
\ldots, x_n$ and writes to address $z$. The operational behaviour of processes
is then determined by the rules of the proof system, for instance cut and
snip allocates a fresh memory cell, spawning a writer process and a reader
process that synchronize through it, and process expressions assigned to right
rules write to memory while those assigned to left rules read from it.

This low-level interpretation of SAX further gives rise, in extensions of SAX,
to type-theoretically justified compiler optimizations such as in-place memory
reuse \cite{}. Extensions of SAX in adjoint logic have also served as the
logical foundation for concurrent language implementations \cite{}. A
mechanization of SAX's proof theory therefore provides higher confidence in the
theoretical foundations underlying these results.


\subsection{Mechanization}
